\subsection{A leitura das duas saídas pelas duas carteiras}
\label{sec:dois-destinos-varrimento}
\index{carteira@carteira!varrimento de saidas@varrimento de saídas}
\index{etiqueta de vista@etiqueta de vista!custo do varrimento@custo do varrimento}

Alice já escolheu as duas saídas do pagamento. A primeira vai para
Bob e a segunda regressa a Alice. As duas carteiras vão ler os
mesmos bytes e chegar a conclusões diferentes. Para acompanharmos
essa leitura, fixemos endereços principais e escolhamos chaves
didácticas pequenas. Os cálculos desta vez fazem-se em Edwards25519,
com as codificações e funções de Monero~\cite{monero-source-v16}.
O valor dos segredos é aqui público para que possamos seguir também
as operações que cada carteira guarda para si.

Façamos \(r=3\). A chave pública da transacção é, pois, \(R=3G\).
Escolhamos para Bob \((k_B^v,k_B^s)=(5,7)\) e para Alice
\((k_A^v,k_A^s)=(11,13)\). Os dois endereços públicos consistem nos
pares de pontos correspondentes. À carteira que prepara a primeira
saída bastam os pontos de Bob. Obtém \(8rK_B^v=8\cdot3(5G)=120G\).
Ao preparar o troco, Alice conhece a sua chave de vista e calcula
\(8k_A^vR=8\cdot11(3G)=264G\). Para o varrimento posterior
ficam, portanto, as derivações
\[
 \mathcal D_B=120G,\qquad \mathcal D_A=264G.
\]
O \(R\) lido em \texttt{extra} será o mesmo nas duas ocasiões.
É a escolha da chave privada de vista que muda a derivação.

Vejamos uma destas contas até aos bytes. A codificação comprimida
de \(120G\), escrita em dois grupos consecutivos de dezasseis bytes,
é
\[
\begin{gathered}
 \texttt{3b1465e5f12a1ce090a208bd8b23b5ae}\\
 \texttt{fc9996783aa2b3a57be82d50a0710708}.
\end{gathered}
\]
A primeira saída tem índice zero. Para formar o seu escalar
acrescenta-se o byte \texttt{00} à sequência acima. Os 32 bytes
de Keccak--256 aplicados a esta entrada são
\[
\begin{gathered}
 \texttt{9bb82456e3dc96ad8b8d3b395f0ee194}\\
 \texttt{e6668eee34fc0f28d4852cd30c95f2a6}.
\end{gathered}
\]
Recordemos que o byte de menor peso se encontra à esquerda.
O último, \texttt{a6}, é o de maior peso. O inteiro que estes bytes
representam está entre \(10l\) e \(11l\). Tirando \(10l\), obtém-se
a codificação canónica de \(h_0\), que se escreve
\[
\begin{gathered}
 \texttt{59718ab4dbfdde3c2c6d8fdbac4b2bc4}\\
 \texttt{e5668eee34fc0f28d4852cd30c95f206}.
\end{gathered}
\]
A operação que se segue depende de quem a faz. A carteira de Alice
soma os \emph{pontos} \(h_0G+K_B^s\) e deixa a chave resultante na
saída. A carteira completa de Bob soma os \emph{escalares}
\(h_0+k_B^s=h_0+7\). Somar sete faz o primeiro byte passar de
\texttt{59} a \texttt{60}, e os restantes conservam-se.
Multiplicando por \(G\) esse escalar, obtemos o ponto produzido por
Alice. A sua codificação é
\[
\begin{gathered}
 \texttt{147b9732507d001dca2badbed0d72910}\\
 \texttt{964558563b988903f6fc59d92330b970}.
\end{gathered}
\]
No pagamento correspondem três papéis diferentes a objectos do
mesmo comprimento. A derivação é a codificação de um ponto comum
às duas partes. O hash reduzido é a codificação de um escalar que
ambas podem refazer. A chave de saída é a codificação de outro
ponto, disponível a todos os nós. É sobre este último que incidirá
mais tarde a autorização da despesa.

Como chega Bob àquela soma se, de começo, ignora sequer se a saída
lhe pertence? A carteira lê \(R\), calcula \(8k_B^vR=120G\) e
forma primeiro a etiqueta. O cálculo acrescenta ao princípio os oito
bytes correspondentes a \texttt{view\_tag}. Ao fim ficam o mesmo
\(120G\) codificado e o índice \texttt{00}. Do resumo toma-se o
primeiro byte tal como sai da função. O quadro reúne os resultados
para cada índice e para cada carteira.
\[
\begin{array}{l|cc}
 &t=0&t=1\\ \hline
 \text{Bob, a partir de }120G &\texttt{92}&\texttt{45}\\
 \text{Alice, a partir de }264G &\texttt{47}&\texttt{05}\\ \hline
 \text{byte presente na saída}&\texttt{92}&\texttt{05}
\end{array}
\]
O \texttt{92} da primeira posição convida Bob a completar a conta.
Subtrai \(h_0G\) ao ponto transmitido e encontra \(7G\), a sua
chave pública de gasto. A segunda posição pede-lhe \texttt{05},
ao passo que a sua derivação e o índice um produzem \texttt{45}.
Bob passa adiante. Alice faz o percurso complementar. Na segunda
posição encontra \texttt{05} e subtrai à chave dessa saída o
múltiplo de \(G\) derivado de \(264G\) com o índice um. O ponto
que lhe fica é \(13G\). É assim que reconhece o troco.

A chave privada de vista permite efectuar todas as comparações do
quadro. A carteira só de vista de Bob chega ao ponto \(7G\) tão bem
como a completa. Para assinar o gasto, a completa disporá ainda do
escalar que, somado a \(h_0\), dá a chave privada da saída.
Distinguem-se assim dois momentos da vida dos 7 de Bob. Ao descobrir
a saída, a carteira compara pontos. Ao gastá-la, usa a abertura que
permite agir sobre esses pontos.

Podemos medir o que a etiqueta poupou. Suponhamos que uma carteira
examina um lote com \(m\) saídas próprias correctamente etiquetadas e
\(n\) saídas alheias. Para cada saída alheia, suponhamos que a
etiqueta esperada pela carteira é uniforme e independente do byte
transmitido. Há 256 valores possíveis e exactamente um coincide.
Por linearidade do valor esperado, o número médio de testes completos
é então
\[
                            m+\frac{n}{256}.
\]
Por exemplo, com duas saídas próprias e 1024 alheias, a média é
\(2+4=6\). Em cada lote concreto o número de coincidências pode
variar. A carteira teve de ler as 1026 etiquetas e de obter as
derivações necessárias para as conferir. A economia incide nas
multiplicações por \(G\) e nas comparações completas que viriam a
seguir. E quando a mesma chave de transacção serve várias saídas
para uma dada chave de vista, a carteira pode calcular a derivação
uma só vez e variar o índice na entrada das funções de hash.

Uma coincidência ocasional de etiqueta tem assim um preço bem
definido. Faz a carteira visitar mais uma chave de saída e subtrair
mais um \(h_tG\). A decisão de a atribuir à conta depende da chave
de gasto encontrada após essa subtracção. Até aqui só tínhamos uma
chave pública de gasto por destinatário. Bob pode desejar uma por
cliente e continuar a reconhecer todos os pagamentos com a mesma
chave privada de vista. Dessa conveniência nascem os sub-endereços.
